Expected shortfall (ES) is not elicitable on its own, but the pair of
value-at-risk (VaR) and ES is \citep{fissler2016higher}, and (VaR, ES)
forecasts are compared through a strictly consistent scoring function for
the pair \citep{nolde2017elicitability}. We work with returns $y$, whose
lower tail is the loss tail, a level $\tau\in(0,1/2)$, and forecasts $(v,e)$
with $e\le v<0$ for the $\tau$-quantile and the $\tau$-expected shortfall,
scored by the FZ0 loss of \citet{patton2019dynamic},
\begin{equation}
 S(v,e;y)=-\frac{1}{\tau e}\ind(y\le v)(v-y)+\frac ve+\log(-e)-1 ,
 \label{eq:fz0}
\end{equation}
the member of the class of \citet{fissler2016higher} whose differences are
invariant to the scale of the returns. Write $\omega=1/(\tau|e|)$ for its
tail slope and $d_t=S(v_{A,t},e_{A,t};y_t)-S(v_{B,t},e_{B,t};y_t)$.

\begin{lemma}[Structure of FZ0 differences]
\label{lem:fz0}
Let $w(v,e)=v/e+\log(-e)$, $c_t=w(v_{A,t},e_{A,t})-w(v_{B,t},e_{B,t})$ and
$D_t=d_t-c_t$. Then
\begin{enumerate}
\item $D_t=\omega_{A,t}(v_{A,t}-y_t)_+-\omega_{B,t}(v_{B,t}-y_t)_+$, which is zero
 unless $y_t\le\max\{v_{A,t},v_{B,t}\}$;
\item if $e_{A,t}\ge e_{B,t}$, so that $A$'s ES is the less conservative,
 then $D_t\ge-\omega_{B,t}(v_{B,t}-v_{A,t})_+$, and if moreover
 $e_{A,t}>e_{B,t}$, then $D_t\to\infty$ as $y_t\to-\infty$; the case
 $e_{A,t}\le e_{B,t}$ is symmetric;
\item $w$ is strictly increasing in $|v|$ and nondecreasing in $|e|$, so
 $c_t\le0$ whenever $A$ is less conservative in both components,
 $v_{A,t}\ge v_{B,t}$ and $e_{A,t}\ge e_{B,t}$, and $c_t<0$ if moreover
 $v_{A,t}>v_{B,t}$.
\end{enumerate}
\end{lemma}

As for intervals and quantiles, the score difference departs from a value
known in advance only on exceedance days (item~1). Unlike them, it is
unbounded on one side whenever the ES forecasts differ (item~2). The same holds for every strictly consistent score in the
characterization of \citet{fissler2016higher}: the difference contains, for
outcomes below both VaRs, a term proportional to $\{G_2(e_A)-G_2(e_B)\}\,y$
with $G_2$ strictly increasing.

\paragraph{What can be certified.}
On days on which $A$'s ES is the less conservative, $d_t$ is bounded below,
so Theorem~\ref{thm:valid} applies to the lower bound; but large positive
deviations inflate its variance term. Flooring the outcome repairs this
without an assumption: replacing an extreme return by a less extreme one
can only lower $d_t$ on these days, so a lower bound computed from the
floored differences remains a valid lower bound for the original ones. Let $\underline y_t\le\min\{v_{A,t},v_{B,t}\}$ be
$\G_t$-measurable, and let $\tilde d_t$ be the score difference evaluated at
$\max\{y_t,\underline y_t\}$, which lies in a predictable bounded range.

\begin{proposition}[One-sided certification]
\label{prop:es-one-sided}
Let $\mathcal L=\{t:e_{A,t}\ge e_{B,t}\}$, a predictable set of days, and
let $\Dbar^{\mathcal L}_n$ be the average of $\Delta_t=\E[d_t\mid\G_t]$ over
the first $n$ days in $\mathcal L$. The lower bound $L_n$ of
Section~\ref{sec:cs}, computed from $\tilde d_t$ on the days in $\mathcal L$,
satisfies $\Pr(\exists n:\Dbar^{\mathcal L}_n<L_n)\le\delta/2$, with no
assumption on the distribution of the returns. The same holds for the upper
bound on the days in the complement of $\mathcal L$.
\end{proposition}

Thus one can show, without assumptions, that a forecaster whose ES is less
conservative was worse on the days on which it was less conservative. The
reverse cannot be shown without an assumption on the loss tail.

\begin{proposition}[No test can certify the less conservative ES]
\label{prop:es-impossible}
Fix forecasts $(v_A,e_A)$ and $(v_B,e_B)$ with $e_A>e_B$, used on every
day, and let the returns be i.i.d. Let $P_1$ be any law with a finite first
moment under which $\E_{P_1}d<0$, so that $A$ is better. Let $\phi$ be any
test based on $y_1,\ldots,y_T$ that rejects the hypothesis
$\E d\ge0$ with probability at most $\delta$ under every law with a finite
first moment for which that hypothesis holds. Then
$P_1(\phi\text{ rejects})\le\delta$.
\end{proposition}

The proof mixes $P_1$ with a small probability $\varepsilon$ of a very
large loss, which reverses the comparison while changing the law of the
sample by at most $T\varepsilon$ in total variation; this is the
obstruction of \citet{bahadur1956nonexistence}. Since the perturbation
changes the probability of an exceedance by at most $\varepsilon$, the
result holds even if the hypothesis is restricted to laws whose exceedance
probabilities for both VaR forecasts lie in any given open neighbourhood of
those under $P_1$, for instance laws under which both VaR forecasts are
nearly calibrated.
Certifying that a less conservative ES forecast is better therefore
requires an assumption on the loss tail. Under a known lower bound on the
returns, flooring the outcome at that bound changes nothing, the score
difference has a predictable range on all days, and the results of
Sections~\ref{sec:cs} and~\ref{sec:rates} apply.

\paragraph{Rates and the Diebold--Mariano test.}
Both carry over. With equal ES forecasts $e_A=e_B=e$ and $v_B=v_A-H$, the
FZ0 difference is $\omega\{-\tau H+\min((v_A-y)_+,H)\}$, the construction of
Theorem~\ref{thm:lower} scaled by $\omega=1/(\tau|e|)$, so separating the forecasters
requires at least $(1-\tau)\{H/(|e|\Delta)\}^2\log\{1/(4\delta)\}$
expected exceedances, where $H/|e|$ is the distance between the VaR
forecasts relative to the ES. On a sample without exceedances the score
difference equals $c_t$, which by item~3 of Lemma~\ref{lem:fz0} favours
the forecaster that is less conservative in both components, strictly when
its VaR is strictly less conservative. The proof of
Proposition~\ref{prop:dm} then applies with $c_t$ in place of the width
difference.
